
annulés par @@M0@@, donc de @@M1@@-Modules, et on a :
@@M2@@
d'où le résultat. On définit donc une fonction sur la catégorie des @@M3@@-Modules cohérents à valeurs dans @@M4@@ en associant à tout @@M5@@-Module cohérent la classe d'un de ses prolongements sur @@M6@@.

Cette fonction est additive. En effet, soit :
@@M7@@
une suite exacte de @@M8@@-Modules cohérents. Soit @@M9@@ un prolongement de @@M10@@ sur @@M11@@; alors il existe un sous-Module @@M12@@ de @@M13@@ prolongeant @@M14@@, et @@M15@@ prolonge @@M16@@. Elle se factorise donc par @@M17@@, ce qui définit l'homomorphisme @@M18@@ cherché. On vérifie immédiatement qu'il est inverse de @@M19@@.

\textbf{Proposition 1.2.} \uline{Soient} @@M20@@ \uline{deux schémas noethériens}, \uline{et} @@M21@@ \uline{un morphisme} plat \uline{de} @@M22@@ \uline{dans} @@M23@@. \uline{On se donne une famille d'éléments} @@M24@@ \uline{de} @@M25@@, \uline{et on considère} \uline{l'homomorphisme} :
@@M26@@
\uline{défini par} :
@@M27@@
\uline{où} @@M28@@ \uline{est l'homomorphisme} «~\uline{image inverse}~» \uline{de} @@M29@@ \uline{dans} @@M30@@, \uline{défini d'après} IV 2.12, \uline{et où} @@M31@@ \uline{est considéré comme} @@M32@@\uline{-module par} IV 2.10. \uline{On suppose que pour tout point} @@M33@@, \uline{les images des} @@M34@@ \uline{dans} @@M35@@ \uline{par} \uline{l'homomorphisme} :
@@M36@@
\uline{engendrent} @@M37@@ \uline{comme groupe abélien}. \uline{Alors} @@M38@@ \uline{est surjective}.

